19
T HE C ONCURRENCY
V IEWPOINT
De finition
Concerns
Models
ProblemsPitfalls
and
Stakeholders
Applicability
Describes the concurrency structure of the system and maps func-
tional elements to concurrency units to clearly identify the parts of
the system that can execute concurrently and how this is coordi-
nated and controlled
Task structure, mapping of functional elements to tasks, interpro-
cess communication, state management, synchronization and integ-
rity, supporting scalability, startup and shutdown, task failure, and
reentrancy
System-level concurrency models and state models
Modeling the wrong concurrency, modeling the concurrency
wrongly, excessive complexity, resource contention, deadlock, and
race conditions
Communicators, developers, testers, and some administrators
All information systems with a number of concurrent threads of
execution
Historically, information systems were designed to operate with little or no
concurrency, running via batch mode on large central computers. However,
a number of factors (including distributed systems, increasing workloads,
and cheap multiprocessor hardware) have combined so that today’s infor-
mation systems often have little or no batch processing and are inherently
concurrent.
In contrast, control systems have always been inherently concurrent and
event-driven, given their need to react to external events in order to perform
333
334
 P A R T III  A V I E W P O I N T C A T A L O G
control operations. It is natural, then, that as information systems become
more concurrent and event-driven, they start to take on a number of charac-
teristics traditionally associated with control systems. In order to deal with
this concurrency, the information systems community has naturally adopted
and adapted proven techniques from the control systems community. Many of
these techniques form the basis of the Concurrency viewpoint.
The Concurrency view is used to describe the system’s concurrency and
state-related structure and constraints. This involves defining the parts of the
system that execute at the same time and how this is controlled (e.g., defining
how the system’s functional elements are packaged into operating system
processes and how the processes coordinate their execution). To do this, you
need to create a process model and a state model: The process model shows
the planned process, thread, and interprocess communication structure; the
state model describes the set of states that runtime elements can be in and the
valid transitions between those states.
Once you have created process and state models, you can use a number of
analysis techniques to ensure that the planned concurrency scheme is sound.
The use of such techniques is typically part of creating a Concurrency view, too.
It’s worth noting that not all information-based systems really benefit
from a Concurrency view. Some information systems have little concurrency.
Others, while exhibiting concurrent behavior, use the facilities of underlying
frameworks and containers (e.g., application servers and databases) to hide
the concurrency model in use.
EXAMPLE Data warehouse systems tend to be batch-loaded overnight
and accessed from a number of desktop machines. These systems do
exhibit concurrent behavior—multiple clients can request information
from the data warehouse concurrently. However, such a system will
typically rely on the underlying database management system to handle
all of the concurrency for it (in any way it chooses). Therefore, the pro-
cess model used is of little architectural significance, and you have little
or no control over it. The interesting aspects of concurrency relate much
more to the design of the physical data model and should be handled
there.
In contrast, however, many of today’s information systems are inher-
ently event-driven, reactive, concurrent systems. This is particularly the
case when considering infrastructure such as middleware products. Systems
of this type typically sit idle until an external event occurs and then process
the event. Given that many external events can occur simultaneously and
C H A P T E R 19  T H E C O N C U R R E N C Y V I E W P O I N T
 335
that the interarrival time of such events may be lower than the time taken to
process them, this kind of information-based system is inherently concur-
rent, with many operations being executed at once.
EXAMPLE Consider an e-commerce system that uses a message-based
approach to processing transaction requests. In such a system, when a
request arrives, it is translated into a message that is queued for the
appropriate functional element that can process it. In order to prevent
message queues from growing too long and to make efficient use of pro-
cessing resources, the processing element will need to process a number
of messages concurrently. In this case, there may be a large number of
concurrent operations within the functional element, each one needing
access to shared resources.
The Concurrency viewpoint is extremely relevant to systems that exhibit
this kind of behavior. Creating a Concurrency view allows the concurrency
design of such systems to be made explicit and helps interested stakeholders
understand concurrency constraints and requirements. It also allows you to
analyze the system to avoid common concurrency problems such as deadlocks
or bottlenecks.
C ONCERNS
Task Structure
The most important aspect of creating a Concurrency view is establishing the
system’s process structure, which identifies the overall strategy for using
concurrency in the system. It defines the set of processes across which the
system’s workload is partitioned and how the functions of the system are dis-
tributed across them. It is also usually necessary to consider the use of
operating system threads within processes or to abstract away from individ-
ual processes and consider groups of similar processes instead.
Note: Throughout this chapter, we use the word task as a generic term to
describe a processing thread—whether it is a single operating system process,
one thread within a multithreaded process, or some other software execution
unit. When the difference is significant, we specifically use the terms process
or thread as appropriate.
336
P A R T III  A V I E W P O I N T C A T A L O G
The aspects of the system’s task structure that this view needs to address
depend very much on the kind of system you are dealing with.
EXAMPLE A complex, small-footprint system may have only one or two
operating system tasks but may need to use a very complex thread
model to meet its efficiency and responsiveness goals. In this case, the
focus of the task structure activity needs to be at the thread level.
A large enterprise system may be composed of literally hundreds of
concurrent processes, many containing dozens of threads. In this sort of
system, the task structure activity needs to be at the level of groups
of similar processes in order to focus on the architecturally significant
aspects of the concurrency.
Mapping of Functional Elements to Tasks
The mapping of functional elements to tasks can have a significant effect on
the performance, efficiency, resilience, reliability, and flexibility of your ar-
chitecture, so this needs careful consideration. The key question to address is
which functional elements need to be isolated from each other (and so placed
in separate processes) and which need to cooperate closely (and so need to
run within the same process).
Interprocess Communication
When functional elements reside within a single operating system process,
communication among them is relatively simple because of their shared
address space. While some coordination may be required (see the Synchroni-
zation and Integrity concern), you can use any number of data structures to
pass information among them. Similarly, a number of easily used control
mechanisms (such as the procedure call and variants of it) can transfer con-
trol among elements as needed.
In contrast, when elements reside in different operating system pro-
cesses, communication among them becomes more complex. This complexity
increases if the processes also reside on different physical machines.
A number of interprocess communication mechanisms can be used to link ele-
ments in different processes, including remote procedure calls, messaging, shared
memory, pipes, queues, and so on. Each has its own strengths, weaknesses, and
constraints, and inappropriate use of these mechanisms can cause problems at the
system level (e.g., message queue latency between processes causing scalability or
throughput problems). In order to deliver a system with an acceptable set of quality
C H A P T E R 19  T H E C O N C U R R E N C Y V I E W P O I N T
 337
properties, the Concurrency view needs to consider and identify the set of interpro-
cess communication mechanisms that will be used to provide the interelement
communication required by the system’s functional structure.
State Management
In many systems, the runtime state of system elements is important to the
correct operation of the system. This is particularly the case for event-driven
systems that exhibit a high degree of concurrency, where business operations
tend to be processed via state machine implementations.
In such systems, a concern of the Concurrency view is to clearly define
the set of states that each functional element of the system can be in at runt-
ime, the set of valid transitions between those states, and the causes and ef-
fects of the interstate transitions. Such careful state management is a major
factor in ensuring reliability and correct behavior for most concurrent sys-
tems. Again, if you are using a formal architectural style, it may define how
the system’s runtime state should be handled.
Note that this concern refers to the state of the runtime elements of the
system (which could be termed the technical state of the system). Another
type of state management important to many information systems is the set
of valid states and transitions for their core persistent information (business
objects—the business state of the system). However, this is a distinct concept
of state, and we refer to persistent object state models as lifecycles to avoid
any confusion between the two. Object lifecycles are discussed in Chapter 18
as part of the Information viewpoint.
Having said this, it is also quite reasonable to consider state management
in the Functional view—after all, the state of functional elements is what
we’re considering. However, our experience is that the design of the system’s
state management usually fits better in the Concurrency view. Those systems
where state is important are usually those where concurrency is important
too, and considering system-level state usually involves the consideration of
the concurrency around it as well.
Synchronization and Integrity
As soon as more than one thread of control exists in the system, it is important
to ensure that concurrent execution cannot result in corruption of information
within the system. This concern applies at a number of levels in the system,
from a shared variable within a multithreaded module at one end of the scale to
critical corporate transaction data in shared data stores at the other.
An important concern for the Concurrency view to address is how concur-
rent activity will be coordinated so that the system operates correctly and
maintains the integrity of the data within it.
338
 P A R T III  A V I E W P O I N T C A T A L O G
Supporting Scalability
In any highly concurrent system, the approach taken to concurrency, synchro-
nization, and state management can have a profound effect on the scalability
that the system can achieve. Too much or too little concurrency can slow a
system down and prevent it from handling heavy workloads efficiently, while
excessive or simply naïve synchronization can result in a system that per-
forms very well for light workloads but grinds to a halt when heavy work-
loads are applied. The problem is that designing high-performance, highly
concurrent systems is difficult and can be an error-prone process. An impor-
tant concern for this view to address is how the concurrency approach used
will support the performance and scalability required while being simple
enough to implement cost-effectively and reliably. We discuss an approach to
achieving this in Chapter 26 on the Performance and Scalability perspective.
Startup and Shutdown
When you have more than one operating system process in your system, star-
tup and shutdown of the system can become more complicated to manage.
Intertask dependencies may mean that tasks need to be started and stopped in
very specific orders so that if some tasks fail to start, others will not be
started. The system startup and shutdown dependencies are an important part
of your concurrency design and need to be clearly understood by developers,
testers, and administrators.
Task Failure
When functional elements reside in different processes or run on different
threads, dealing with element failure can be complex. This is because an
element in one task cannot rely on another task being available when it
needs to communicate with it, whereas when an element calls another one
in the same task, it knows it will be there. Your concurrency design needs
to take into account this added possibility of failure and ensure that the
failure of one task doesn’t bring the entire system to a halt. In order to
address this concern, you need a system-wide strategy for recognizing and
recovering from task failure.
Reentrancy
Reentrancy refers to the ability of a software element to operate correctly
when used concurrently by more than one processing thread. This is primarily
a concern for software developers when designing their software elements.
From an architectural perspective, reentrancy is an important constraint for
C H A P T E R 19  T H E C O N C U R R E N C Y V I E W P O I N T
 339
certain elements, so the architecture must clearly define which modules need
to be reentrant and which do not.
EXAMPLE If you are developing an e-mail server, the ability to support
a great deal of concurrency is likely to be a key concern. Without this, it
will be hard to use the e-mail server for large user populations who will
want to send and receive e-mail simultaneously. You can take a number
of approaches to achieve such concurrency, but for the sake of argu-
ment, let’s assume that you have decided to implement the server by
using a single operating system process and many (perhaps hundreds
of) concurrent operating system threads running within it: some send-
ing e-mail, some receiving e-mail, and some managing the server’s
internal state.
In this sort of environment, it is crucial to decide which of the ele-
ments of your system have to be reentrant and which do not. Any ele-
ment involved in sending and receiving e-mail (e.g., a name resolution
library that translates e-mail domains to network addresses) will need to
be reentrant to ensure that it can be used simultaneously by many send-
ing and receiving threads. Without such a guarantee, the name resolu-
tion library could be the source of many subtle problems later if its
internal state could be corrupted by concurrent access.
The reentrancy needs of your architecture can also affect which third-party
software elements you can use within the system and where you can use them.
Stakeholder Concerns
Typical stakeholder concerns for the Concurrency viewpoint include those
shown in Table 19–1.
TABLE 19–1 S TAKEHOLDER C ONCERNS FOR THE C ONCURRENCY V IEWPOINT
Stakeholder Class
 Concerns
Administrators
 Task structure, startup and shutdown, and task failure
Communicators
 Task structure, startup and shutdown, and task failure
Developers
 All concerns
Testers
 Task structure, mapping of functional elements to tasks, startup
and shutdown, task failure, and reentrancy
340
 P A R T III  A V I E W P O I N T C A T A L O G
M ODELS
System-Level Concurrency Models
The Concurrency view maps the functional elements onto runtime execution
entities via a concurrency model. The concurrency model typically contains
the following items.
Processes: In this context, the term process refers to an operating system
process, that is, an address space that provides an execution environ-
ment for one or more independent threads of execution. The process is
the basic unit of concurrency in the design of the system. At the architec-
ture level, the processes are normally assumed to be isolated from each
other so that if one process wants to affect the execution of another, it
must use an interprocess communication mechanism.
Process groups: At the architecture level, it can often be useful to group in-
dividual processes so that a collection of closely related processes can be
considered as a single entity at the system level. This can provide a useful
abstraction that allows less important concurrency concerns to be deferred
until subsystem design. An example is a database management system
(DBMS). The important point from the system level is that the DBMS is a
functional unit, accessed via well-defined interfaces, that runs in its own
process or group of processes. However, the details of the exact number of
processes it uses (e.g., how many logging processes run within the DBMS)
and the function of each are almost certainly irrelevant to the architec-
ture—indeed, this will probably be decided by a technical specialist later in
the design process. Using a process group in this situation makes it clear
that a group of related processes will be used but defers the details of the
set until later. The other common use for process groups is simply as a hi-
erarchical structuring technique for large or complex systems that contain
many processes. All of the processes may need description, but the use of
process groups can make the process model easier to comprehend.
Threads: In this context, the term thread refers to an operating system
thread, that is, a thread of execution that can be independently scheduled
within an operating system process. Threads are known as lightweight
processes by some operating systems. At the level of system architecture,
threads can often be ignored, with the details of their use being the
responsibility of subsystem designers (perhaps with you guiding their
use via design patterns in the Development view). However, for some
systems you do want to model the use of threads in at least some parts of
the system. Threads are normally represented in process models via a
decomposition of a process.
C H A P T E R 19  T H E C O N C U R R E N C Y V I E W P O I N T
 341
Interprocess communication: When processes are running, they are
assumed to be isolated from each other so that one process cannot
change anything in another process. However, in most concurrent sys-
tems, processes do need to interact in order to coordinate their execution,
request services from each other, and pass information among them-
selves. They achieve these interactions via a number of interprocess
communication mechanisms (“IPC mechanisms”), which are the connec-
tors in the system’s runtime architecture.
The mechanisms available vary depending on the underlying tech-
nology platforms in use. However, interprocess communication mecha-
nisms generally fall into one of these groups.
• Procedure call mechanisms are all some variation on an interprocess
function call and are usually based on some form of remote procedure
call or some sort of message-passing operation.
• Execution coordination mechanisms allow two or more processes (or
threads) to signal to each other when certain events occur.
Coordination mechanisms include semaphores and mutexes and are
typically limited to coordination between processes or threads running
on the same physical machine.
• Data-sharing mechanisms allow a number of processes to share one or
more data structures and access them concurrently (possibly
coordinating this access via coordination mechanisms). Data-sharing
mechanisms include shared memory, distributed tuple spaces (such as
Linda and more recent implementations such as GigaSpaces), and
simple, traditional mechanisms such as client/server databases and
shared file storage.
• Messaging mechanisms are related to data-sharing mechanisms, but
rather than placing data structures in a shared space for concurrent
access, they transmit data structures from one task to another.
Messaging systems normally implement one or both of two well-
defined messaging models: queuing and publish/subscribe. Queuing
introduces a “first in, first out” queue structure between producers
and consumers where consumers destructively read messages from
the queue (i.e., a message is delivered to only one consumer).
Publish/subscribe introduces a “topic” or “bus” between producers
and consumers where the consumers indicate the types of messages
that they are interested in and a message is consumed by all
consumers interested in it.
As the architect, you need to choose the interprocess communication
mechanisms carefully because of the impact that they can have on the quality
properties that the system exhibits (such as its performance, scalability, and
reliability). The IPC mechanisms also impose significant constraints on the
functional elements that use them, and so it is important to choose them early
so that these constraints can be taken into account.
342
 P A R T III  A V I E W P O I N T C A T A L O G
NOTATION You can represent the Concurrency view in a number of ways. Some
of the more common notational approaches include UML and other formal nota-
tions, along with less formal notations, which we describe briefly here.
UML: UML’s concurrency modeling facilities are rather simple but do
include the notion of an active object (i.e., an object with a thread of con-
trol). There are a number of ways that concurrency structures can be rep-
resented in UML, including stereotyped packages, components, and
classes, but unfortunately no approach has become a standard. W e have
used a number of conventions in our models and have found that stereo-
typed active components are a good place to start when modeling pro-
cesses and threads. We sometimes also add a process group stereotype to
represent a group of related processes (such as a database engine) if this
is a useful abstraction for the system we are working on.
Simple examples of interprocess communication, like remote proce-
dure calls, can be represented by using standard UML intercomponent
associations, with arrowheads indicating the direction of communication
(and possibly using tagged values on the association to make the com-
munication mechanism clear). More complex forms of interprocess com-
munication (shared memory, semaphores, and so on) can be represented
quite effectively by introducing further stereotypes and showing associa-
tions between the components in the tasks and the interprocess commu-
nication mechanisms they use.
Figure 19–1 shows an example of UML being used for a concurrency
model.
This model shows how the system is implemented by using three
processes (a client, a statistics service, and a statistics calculator) along
with a process group to implement the Oracle DBMS instance. The con-
current activity between the Statistics Accessor and Statistics Calculator
components needs to be coordinated because they are in different pro-
cesses; a mutex is used to achieve this. The illustrated scenario is very
simple, and there is little or no architecturally significant thread design
in this model. Figure 19–2 shows a more involved model with more ar-
chitecturally significant threading.
The concurrency model shown in Figure 19–2 illustrates a case
where the process structure is very simple, namely, two processes that
communicate via a socket stream. However, the thread structure in the
DBMS Process instance is architecturally significant, and its structure
and interthread coordination strategy need to be documented and
explained. The model shows that there is a single thread containing the
Network Listener component, which communicates with between 1 and
40 threads that contain the four main query-processing components via
an interprocess communication queue. The Disk I/O Manager component
FIGURE 19–1 C ONCURRENCY MODEL D OCUMENTED BY U SING UML
is hosted on its own thread, and there may be up to 10 instances of this
running. The Data Access Engine component communicates with the
Disk I/O Manager instances via the shared memory mechanism.
Formal notations: The real-time and control systems research community has
created a number of concurrency modeling languages that allow the creation
and analysis of process models. A number of these languages, such as LOTOS,
Communicating Sequential Processes (CSP), and the Calculus of Communicat-
ing Systems (CCS), are formal and represented textually. Most of these lan-
guages are mathematical and fairly abstract, and they aren’t widely used in
information systems development. While this doesn’t mean they can’t be use-
ful, we have yet to come across a large-scale industrial application of them to
information systems. The problem with using these languages in practice is
often the need to teach them to the interested stakeholders, and there is always
a need to ensure that the representation and analysis that they allow will be
useful for the specific situation to which you are trying to apply them.
Informal notations: In our experience, by far the most common notation
used to represent process models is an informal one, created by the
author of the model. Given the relatively small number of object types in
a process model, an informal notation invented for the problem at hand
often works very effectively as a communication tool as long as it is
explained clearly. The notation needs to capture processes, process
groups, threads, and the interprocess communication mechanisms in
use. As long as the notation to represent each of these concepts is well
defined, an informal notation often has much to recommend it. In partic-
ular, the notation can be kept simple and avoids the potentially awkward
process of bending a general-purpose notation like UML to represent the
model being described. The risk with informal notations is that they are
never clearly defined and so lead to ambiguous descriptions.
ACTIVITIES
Map the Elements to the Tasks. The first step when creating your process
model is to work out how many processes you need and to decide which func-
tional elements will run in which processes. In some cases, this is a straight-
forward job—each functional element ends up being a process (or perhaps a
process group), or all of the elements end up in a single process. In other
C H A P T E R 19  T H E C O N C U R R E N C Y V I E W P O I N T
 345
cases, there is a complex N:M mapping between functional elements and pro-
cesses, with some elements partitioned between processes and other elements
running in shared processes. The important point about this mapping is that
you should introduce concurrency only where it is actually required. Concur-
rency adds complexity to the system and adds significant overhead to inter-
element communication when it must cross process boundaries. Therefore,
add more processes to your system only if you need them for distribution,
scalability, isolation, or other reasons guided by the requirements for your
system.
Determine the Threading Design. The term threading design refers to the
process of deciding on the number of threads to include in each system pro-
cess and how those threads are to be allocated and used. In most cases,
threading design is not something that the architect needs to get directly
involved in—this is usually the job of the subsystem designers. However, you
may get involved in designing and specifying general threading approaches or
patterns that should be used at various points in the system in order to meet
the system’s required quality properties or to ensure consistency across the
implementation.
Define Mechanisms for Resource Sharing. As soon as concurrency is
introduced into the system, you must carefully consider how to share
resources between concurrent threads of execution. Resource sharing is
considered in some other parts of the architecture, too (notably, the Infor-
mation view), and the two activities might be best tackled as a single task.
This isn’t a book on concurrent computing, so we don’t have space to dis-
cuss all of the options and potential pitfalls to consider when sharing
resources. The simplistic advice is simply that whenever a resource (such as
a piece of data in memory, a file, a database object, or a piece of shared
memory) is shared among two or more concurrent threads of execution, it
must be protected from corruption. This is usually achieved with some form
of locking protocol. As with threading design, the details of resource sharing
are rarely architecturally significant. Your role in relation to this is to ensure
that suitable resource-sharing approaches are used where necessary and
that the approach used is suitable in the overall context of the system and
does not produce unacceptable side effects for the system as a whole.
Define the IPC Mechanisms to Use. In most concurrent systems the tasks
need to communicate frequently, and so, along with deciding how to share re-
sources between tasks, you will need to consider what communication is
needed between them and which interprocess communication mechanisms
you will use to enable it. Again, we don’t have space here to discuss all of the
options and the tradeoffs that they imply, but as is usually the case, a simple
and regular scheme that minimizes the amount of intertask communication is
likely to be the best choice. Better still is using a library or framework (such
346
 P A R T III  A V I E W P O I N T C A T A L O G
as an implementation of the Actor pattern) to avoid having to deal with a lot
of this complexity yourself. While this sounds like simplistic advice, imple-
menting complex intertask communication correctly is a very difficult job best
left to specialists. As with other architectural concerns, your focus needs to be
on defining a common system-wide approach and on reducing the risk in-
volved in its implementation.
Assign Priorities to Threads and Processes. Some tasks in your system may
be more important than others. If you have tasks of different importance run-
ning on one machine, you need to control their execution so that the more
important work gets done before the less important work. The normal method
for achieving this is to use the facilities of the underlying operating system to
assign priority levels to the different threads and processes. All modern operat-
ing systems provide this feature in roughly the same way. Tasks are explicitly
or implicitly given runtime priorities. When the operating system’s thread
scheduler is choosing tasks to run, it considers the higher-priority tasks before
the lower-priority ones, thus getting the important work done first. If you can
avoid assigning explicit priorities to threads, in general do so—processing pri-
orities can add a lot of complexity to your process model and can introduce sub-
tle but serious problems such as priority inversion. However, sometimes you
can’t avoid it. In these cases, keep the assignment of priorities as simple and as
regular as possible, and analyze and prototype your approach to make sure that
you aren’t introducing problems worse than the one you’re trying to solve.
Analyze Deadlocks. Having introduced concurrency into the system, you
have also introduced the risk of the entire system grinding to a halt in unex-
pected ways. Whenever you have concurrency in the presence of shared
resources, you always have the possibility of deadlock. You can use a number
of modeling and analysis techniques to try to spot potential deadlocks. An
example of such a technique is Petri Net Analysis, which allows you to create
a model of your processing threads and shared resources and then analyze
the model to catch potential deadlock situations. With experience, it is also
usually possible to perform effective deadlock analysis through careful, infor-
mal consideration of your concurrency model.
Analyze Contention. Whenever you have a number of tasks and shared
resources, you almost always find contention. Contention occurs between
tasks when more than one task requires a shared resource concurrently. The
introduction of coordination mechanisms (such as mutexes) inevitably intro-
duces contention when workloads are high. If contention rises beyond a cer-
tain point, the system will slow dramatically, and little useful work will get
done. In order to avoid this during normal operation, you need to analyze
your shared resources from this point of view. The basis of the technique is to
identify each of your possible contention points. Then, for each, estimate the
likely number of concurrent tasks contending for the resource and for how
C H A P T E R 19  T H E C O N C U R R E N C Y V I E W P O I N T
 347
long each will need the resource. This allows you to establish the likely wait
times that each task experiences at each point and then to estimate how such
contention will affect your processing times and throughput. Repeating the
exercise for different workloads allows you to estimate the maximum theoret-
ical workloads your system can possibly support. Once you understand the
potential for the system becoming overloaded, you can design mechanisms
into your software to handle such conditions gracefully (such as implement-
ing the Circuit Breaker pattern).
State Models
A state model is used to describe the set of states that a system’s runtime elements
can be in and the valid transitions between the states. The set of states and transi-
tions for one runtime element is known as a state machine, and the collection of all
of the interesting state machines for your system forms the overall state model.
Usually, you will find that each system task identified in the concurrency
model will have one or at most a few functional units mapped to it that are
effectively in control of the task. These functional units normally have the
system’s interesting state models associated with them. If you create a state
model, be sure to focus on these system elements so that the state model
describes only architecturally significant information. You don’t need to cap-
ture all of the state machines inside all of the system’s elements; the AD
needs to describe only state that is visible at the system level, not state that is
hidden inside the system’s elements.
An important decision to make before you start creating the state model is
the set of semantics you want to use in your state machines. Modern state
modeling notations (in particular, UML’s statechart, discussed later) allow
you to introduce a mind-boggling degree of complexity. You need to use such
notations carefully if you want to produce a comprehensible model.
A basic state machine in the state model would normally contain the fol-
lowing types of entities.
State: A state is an identifiable, named, stable condition during a runtime
functional element’s lifetime. States are normally associated with waiting
for something (an event) to occur or performing some sort of operation.
Transition: A state transition defines an allowable change in state, from
one state to another, following the occurrence of an event. From a model-
ing point of view, transitions are normally considered to occur in zero
time and so cannot be interrupted.
Event: An event is an indication that something of interest has happened
in the system (and is normally recognized by an operation being invoked
on an element or a time period ending). Events are the triggers that
cause transitions between states to occur.
348
 P A R T III  A V I E W P O I N T C A T A L O G
Actions: Actions are atomic (noninterruptible) pieces of processing that
can be associated with a transition (so an event causes the transition to
occur, and then an action is executed as part of the state transition).
More sophisticated state modeling notations allow additional modeling
elements such as guards (Boolean conditions governing state transitions),
activities (long, interruptible items of processing that can be associated with
states), and hierarchical states.
NOTATION State models are typically represented by a graphical notation
derived in some way from the classic state transition diagram. The most popular
variant in use today is probably the UML notation for representing state, the
statechart. At the end of this subsection we briefly discuss other graphical nota-
tions as well as some nongraphical ones, but first we focus on UML’s statechart.
 UML: A statechart is a flexible notation that can be used in a number of
different ways at differing levels of sophistication. Deciding which parts
of the notation to use is an important step before getting too far into the
modeling process. Figure 19–3 shows a UML statechart that represents
the state model for a calculation engine.
This statechart shows much of the important notation for a UML
statechart, with a composite primary state, concurrent state management,
and the use of start and end pseudo states to indicate how the element’s
lifecycle begins and ends. The single top-level state (Running) is entered
when the element is started and exited when a shutdown event is
received (the reset() action is performed as part of that transition).
The Running state has been decomposed into four substates that
comprise the business of running this element: Waiting for Data, Cali-
brating Metrics, Calculating, and Distributing Results. The transition
arrows indicate the possible transitions between states (along with the
events that cause the transitions and the actions that will be executed).
The Calculating state is interesting because it is a concurrent state, as
you can see from the dashed line that bisects it. This means that while in
the Calculating state, the element is actually in two concurrent substates
(Calculating Values and Calculating Risk). When the activity associated
with these states completes, the transition from the states is taken, and
when both are complete, the element can leave the Calculating state.
There are two architecturally significant aspects to this state machine.
• If new input data becomes ready when the element is in the
Calculation state, in-progress results are discarded (by executing the
reset() action), and calculation starts again. In contrast, if this
occurs while results are being distributed, the distribution process is
not interrupted.
FIGURE 19–3 EXAMPLE S TATECHART FOR A CALCULATION E NGINE
• No matter the state of the element, if a shutdown event is received, all
processing immediately stops, the state is reset, and the element exits.
Of course, whether or not these facts are architecturally significant
depends on the situation. However, we can make a reasonable argument
that these facts are visible at the system level and thus can affect or be
relied on by other system elements, and therefore, these facts need to be
captured as part of the architecture.
An interesting point to note about the UML statechart is that its abil-
ity to show hierarchical state composition allows you to express architec-
tural constraints on state models without needing to define the entire
model. The statechart in Figure 19–4 illustrates this point.
This statechart distills one of the architecturally significant features
from the statechart in Figure 19–3, namely, that a shutdown event must
be immediately responded to in any running state and a reset of the ele-
ment performed as part of shutdown. In effect, this documents an architec-
tural constraint that the designer of the corresponding part of the system
350
 P A R T III  A V I E W P O I N T C A T A L O G
Running
shutdown/reset()
FIGURE 19–4 A RCHITECTURAL C ONSTRAINT S TATECHART
must respect; but while clearly defining this constraint, the statechart
leaves the details of the lower-level states to the subsystem designer.
Other graphical notations: In addition to UML, many other graphical
notations exist for modeling state. Some of the better-known ones
include simple state transition diagrams, Petri Nets, SDL, and David
Harel’s original Statecharts. All of these notations have strengths and
weaknesses when compared to each other and to UML statecharts, and
they are worth considering if UML statecharts cause you problems. How-
ever, the standardization of and wide familiarity with UML statecharts
means that in general they should probably be your first choice. The Fur-
ther Reading section at the end of this chapter contains some references
to information about these alternative notations.
Nongraphical notations: In principle, all of the graphical notations can be
represented in a textual form (and indeed many graphical state modeling
notations do define an equivalent textual form). Similarly, a number of
primarily textual formalisms for modeling and analyzing state can be
represented in a graphical form (for an accessible example, look at the
Finite State Processes language). A textual state model can be useful
when the model needs to be processed in some way by machine, but for
human readers it is almost always better to use a graphical notation
where possible.
ACTIVITIES
Define the Notation. Before starting to create your state model, spend some
time working out your needs for the modeling notation and defining how you
will use it.
Identify the States. The primary activity when creating a state model is to
work out what states your system elements can be in and the processing (if
any) associated with each state. Beware of accidentally modeling activities as
P ROBLEMSC H A P T E R 19  T H E C O N C U R R E N C Y V I E W P O I N T
 351
states; this is a common modeling mistake. If in doubt, try considering your
state machine as a UML activity diagram. If you can do this, you have proba-
bly modeled activities rather than states. When performing state identification
at the architectural level, focus on the states that are visible from outside the
element and thus have a system-wide effect.
Design the State Transitions. Once you know what states your elements
can be in, design a set of transitions that allows them to move between the
states correctly. For each transition, clearly identify how the transition is
triggered and any (atomic) actions that must be performed as a side effect of
traversing it. Make sure that the events and actions you identify can be sup-
ported by the operations and state of the element for which you are
designing a state machine.
AND P ITFALLS
Modeling the Wrong Concurrency
When considering the concurrency design of a system, it is easy to get
bogged down in the details of the internal concurrency and state design of
each element. It’s not part of your job as an architect to design detailed
thread models that define how individual threads in a server will be allo-
cated, used, and freed, along with all of the coordination between them.
Remember that your role is to concentrate on the system as a whole rather
than all of the details of each element. The concurrency with which you
should be concerned is the architecturally signi ficant concurrency, that is,
the overall concurrency structure, the mapping of functional elements to that
structure, and the system-level state model. You may also be involved in
specifying common approaches such as design patterns that need to be used
for the concurrency within elements, but in general you should not need to
design all of the details—this will only distract you from the system-level
problems (which are often quite enough to worry about).
RISK REDUCTION
 Focus on architecturally significant aspects of concurrency.
 Involve the lead software developers as early as possible so they can
work on the more detailed aspects of this problem.
Modeling the Concurrency Wrongly
Meaningful concurrency models can be quite difficult to create, so it is impor-
tant to spend the time and effort required to create a good model. To be use-
352
 P A R T III  A V I E W P O I N T C A T A L O G
ful, the models you create should use your notation correctly and be a valid
representation of the situation you are representing. The following are some
of the common modeling mistakes to be aware of.
Modeling activities from your system as states in your state model and so
accidentally creating an activity diagram rather than a state model. This
is a common confusion when people are new to state modeling.
Having a model with nonterminal states that can’t be exited because the
required events will never occur. This indicates that some conditions
haven’t been thought through.
Having transitions that cannot be traversed due to an invalid combina-
tion of events and conditions. This can indicate a misunderstanding or
that something is missing.
The existence of a large number of transitions with just trigger events or
just actions. Such transitions are sometimes required, but having a large
number of them often indicates that there are too many states in the
state model.
RISK REDUCTION
 Before you start, take the time to understand your concurrency modeling
notation (and the way that your tool implements the notation, if you’re
using a modeling tool).
Watch out for states, transitions, or actions that are difficult to name or
that seem to need the wrong sort of name (such as a verb for a state).
These suggest that there is something wrong with the model.
Walk through your models, “playing computer” in order to validate them
and check for missing or incorrect elements. If your tool offers model
animation facilities, these are a more reliable way of achieving this.
Excessive Complexity
Simplicity should always be an aim when designing a system. Simple designs
are easier to create, analyze, build, deliver, and support. However, this is par-
ticularly important when considering concurrency because it is fundamentally
difficult to understand. As we have seen, the price of complex concurrency
can be very high at design time, implementation time, and beyond. More soft-
ware engineering hours have probably been wasted on reworking problematic
concurrency than on almost anything else. Simplicity in your concurrency
approach will have a major positive impact on the amount of effort required to
deliver and support your system.
C H A P T E R 19  T H E C O N C U R R E N C Y V I E W P O I N T
 353
RISK REDUCTION
 Be sure that all of the concurrency you introduce is justified in terms of
stakeholder benefits.
 When designing state models, use the simplest subset of notation possible
to capture your state machines in order to encourage a simple state model.
Resource Contention
Resource contention usually manifests itself as long wait times for shared
resources or excessive activity in small, specific parts of the system (colloqui-
ally known as hot spots). Careful and early analysis of the concurrency model
for potential contention can help you avoid such problems, but in reality, as
soon as one resource contention point is eliminated, the next one will emerge.
Therefore, tackling resource contention is normally a process of reducing the
contention to an acceptable level.
RISKREDUCTION
Analyze your system as it is being designed to spot resource contention
as early as possible, and design around it. Use your usage scenarios to
predict which parts of the system are likely to encounter high levels of
concurrency, and focus your attention in these regions.
Reduce contention by decomposing locks on large resources into a number
of finer-grained locks, thus reducing the amount of time locks are held.
Consider alternative locking techniques such as optimistic locking that
reduce the time locks are held.
Eliminate shared resources where you can, or consider making them
immutable to avoid the need to lock them when accessed by multiple
tasks.
If possible, reduce the amount of concurrency you need around problem-
atic contention points.
Consider whether it is possible to avoid locking by using approximately cor-
rect results that may be in the process of being updated or whether it would
be possible to allow copies of data to be loosely rather than tightly consis-
tent, to avoid locking them all simultaneously during the update process.
Deadlock
Deadlock occurs when one task, A, requires access to a resource that has
already been locked by another task, B, and task B is also waiting for access
to a resource that has been locked by task A. We describe these two tasks as
354
 P A R T III  A V I E W P O I N T C A T A L O G
“deadlocked,” and neither of them will make progress until the deadlock is
broken by one of the tasks releasing one of the locks (possibly by the task
being terminated by a supervisor task). As with resource contention, you can
often avoid deadlock through early and thorough analysis of the system.
Danger points are those parts of the system where different types of process-
ing tasks need access to a number of the same resources. When you find
potential deadlock points, you will probably need to redesign the system to
avoid the problem.
RISK REDUCTION
 Where possible, ensure that resources are always allocated to tasks and
locked in a fixed order.
 Attempt to isolate parallel tasks in such a way that deadlock between
them is impossible.
 Reduce the number, scope, and duration of locks held where possible
(or even use immutable data structures if possible, to completely avoid
locking).
 Certain commercial products that use locks (such as database manage-
ment systems) provide significant assistance with handling deadlock—in
most cases, recognizing it and breaking it by terminating one or more of
the problematic transactions. These technologies can be very useful
when dealing with deadlock, but their use often needs to be carefully
designed into the system so that such deadlock recovery actions are han-
dled correctly.
Race Conditions
A race condition is problematic behavior that results from unexpected depen-
dence on the relative timing of events. It usually occurs when two or more
tasks are attempting to perform the same action concurrently. The tasks race
for the resource, and the first one to reach the appropriate point in the pro-
gram code wins and performs the action.
Race conditions are problematic only when they are unplanned because the
system has not been designed to cope with more than one task performing the
action concurrently. In these cases, information can be corrupted or lost, and
the system can behave in unpredictable ways. A classic example is a system-
wide data structure in an operating system process that a number of threads
can update. If multiple tasks try to update the data structure concurrently
(e.g., to increment a counter indicating the number of requests accepted), the
resulting value will be undefined and very likely incorrect.
C H A P T E R 19  T H E C O N C U R R E N C Y V I E W P O I N T
 355
RISK REDUCTION
 Ensure that there are no unprotected, shared system-level resources that
can cause race conditions.
 Use immutable data structures where possible to avoid the possibility of
race conditions.
 Automatically introduce protection mechanisms for all potentially shared
resources.
 Ensure that the definition of each element interface clearly states
whether or not the interface is reentrant.
C HECKLIST
Is there a clear system-level concurrency model?
Are your models at the right level of abstraction? Have you focused on
the architecturally significant aspects?
Can you simplify your concurrency design?
Do all interested parties understand the overall concurrency strategy?
Have you mapped all functional elements to a process (and thread if
necessary)?
Do you have a state model for at least one functional element in each
process and thread? If not, are you sure the processes and threads will
interact safely?
Have you defined a suitable set of interprocess communication mecha-
nisms to support the interelement interactions defined in the Functional
view?
Are all shared resources protected from corruption?
Have you minimized the intertask communication and synchronization
required?
Do you have any resource hot spots in your system? If so, have you esti-
mated the likely throughput, and is it high enough? Do you know how
you would reduce contention at these points if forced to later?
Can the system possibly deadlock? If so, do you have a strategy for rec-
ognizing and dealing with this when it occurs?
FURTHER R EADING
The area of concurrency has been studied and written about widely, although
not many books consider it from an architect’s perspective.
356
 P A R T III  A V I E W P O I N T C A T A L O G
A good overview of concurrency (albeit with a Java-specific slant) and a
good introduction to (fairly formal) modeling and analysis appear in Magee
and Kramer [MAGE06]; this book also introduces the Finite State Processes
language mentioned earlier. Unfortunately, Cook and Daniels [COOK94] is
out of print; however, it has recently made a welcome reappearance as a freely
available online book at www.syntropy.co.uk/syntropy, as it contains a lot of
good advice on modeling and a particularly good discussion of using state-
charts to model object state. You can find a lot of good UML-specific advice
about state modeling in Rumbaugh et al. [RUMB99], which is organized as a
reference so it’s easy to find definitions of the various UML elements
involved.
Each of the visual formalisms has its own following and literature.
Girauld and Valk [GIRA02] is a relatively academic text that explains how to
apply Petri Nets to the analysis of concurrency characteristics, while the SDL
Forum Society Web site [SDL02] is a good starting point for finding out more
about SDL. A fairly recent reference on CSP is Roscoe [ROSC97]; the defini-
tive book on CCS is still Milner [MILN89]; and the original reference for state-
charts was Harel’s paper [HARE87].
We reference it elsewhere too, but Michael Nygard’s book [NYGA07] has
much to say on the topic of safely introducing concurrency into systems. A
rich collection of design patterns for creating concurrent systems is docu-
mented in [SCHM00], and a thorough programming-level “nuts and bolts”
introduction to concurrency practice can be found in [BRES09]